\chapter*{Цели и задачи}
\addcontentsline{toc}{chapter}{Введение}
\label{ch:intro_lab5}

\textbf{Цель лабораторной работы:} изучение теоретических принципов и инструментальных средств для построения пайплайна предварительной обработки данных.

\textbf{Основные задачи:}
\begin{itemize}
    \item предварительная обработка данных;
    \item изучение библиотек для предварительной обработки данных;
    \item масштабирование признаков;
    \item представление категориальных данных;
    \item построение пайплайна для предварительной обработки данных.
\end{itemize}

В качестве набора данных для выполнения индивидуального задания выбран набор данных \textbf{Abalone} (морское ушко), который содержит информацию о возрасте моллюсков, определяемом по количеству колец на раковине. Данный набор содержит как числовые, так и категориальные признаки, что позволяет продемонстрировать полный спектр операций предварительной обработки данных.

\endinput